\ssr{РЕФЕРАТ}

Расчётно-пояснительная записка объёмом 44 с., в том числе 14 рис., 4 табл., 20 источников, 2 прил.

МОДЕЛИРОВАНИЕ ОСВЕЩЕНИЯ, ТРАССИРОВКА ЛУЧЕЙ, ПРЕЛОМЛЕНИЕ СВЕТА, КАРТЫ ОСВЕЩЕННОСТИ, КАУСТИКИ, ЗАКОН СНЕЛЛИУСА, АППРОКСИМАЦИЯ ФРЕНЕЛЯ, МОДЕЛЬ ЛАМБЕРТА.

Разработана программная система для физически корректного моделирования световых паттернов гирлянд с учетом двукратного преломления в защитных колпачках (сферический слой из двух концентрических сфер). Реализован программный комплекс на языках C++ и Python с использованием pybind11, обеспечивающий трассировку лучей с учетом преломления (закон Снеллиуса, аппроксимация Френеля), формирование карт освещенности для диффузных поверхностей (модель Ламберта), поддержку стохастического рассеивания. Графический интерфейс на Tkinter предоставляет инструменты создания паттернов источников света и интерактивное управление камерой.

В исследовательской части показано, что вычислительная сложность алгоритма составляет $O(N \cdot M)$, где $N$ --- число источников, $M$ --- число объектов. Время расчета линейно зависит от количества источников (0.36 сек/источник) и объектов (0.35 сек/объект), слабо зависит от разрешения карт (прирост 11.5\% при увеличении разрешения в 5 раз). Лимитирующим фактором производительности является этап генерации лучей (80\% времени), в то время как запись в текстурные атласы вносит менее 12\% задержки.

Областью применения являются предварительная визуализация декоративного освещения в архитектурном дизайне, планирование световых инсталляций и образовательные приложения для демонстрации принципов геометрической оптики.